\begin{helper}
Let \(V\) be finite and let \(a,b,c\in V\) be pairwise distinct.  The map
\[
q\in\mathbb R^V\longmapsto
\bigl(q_a,q_b,q_c,(q_v)_{v\in V\setminus\{a,b,c\}}\bigr)
\in \mathbb R^3\times\mathbb R^{V\setminus\{a,b,c\}}
\]
preserves finite product Lebesgue measure.
\end{helper}
\begin{proof}
Because \(a,b,c\) are pairwise distinct, the list consisting first of
\(a,b,c\) and then all elements of \(V\setminus\{a,b,c\}\) is a reindexing of
the finite set \(V\).  The displayed map is exactly the coordinate
permutation associated to that reindexing, written with the first three
coordinates separated from the remaining coordinates.  Finite product
Lebesgue measure is invariant under coordinate permutations, and the standard
product identification between
\(\mathbb R^{\{a,b,c\}\sqcup (V\setminus\{a,b,c\})}\) and
\(\mathbb R^3\times\mathbb R^{V\setminus\{a,b,c\}}\) is also
measure-preserving.  Composing these two measure-preserving identifications
gives the stated measure-preserving map.
\end{proof}
